\section{Discussion and Limitations}

The proposal-policy comparisons do not establish a reliable held-out
improvement from parent conditioning or the additional prompt feedback.
EVO-P already uses fitness to select its archive; EVO-F adds numerical
summaries and case descriptors, after which parent histories can diverge.
The study therefore compares complete search policies, rather than
isolating the causal effect of numerical values in otherwise matched
prompts.

The frozen-program evidence is more specific about behavior than
discovery. Scoring changes often vanish before the objective is
measured; when they survive, rarity's functional form changes how
coverage is valued. This connects transfer to the representation,
but leaves archive conditioning and the training corpus as competing
explanations for which programs a policy selects.

The scope is limited by five runs per strategy, one shared training
corpus and a restricted DSL. Matched proposal attempts yield unequal
numbers of unique executable candidates. External evaluation covers
24 topologies under two weight regimes, with few certified optima;
T1/T2 change both cost magnitudes and their relationship to degree.
These constraints limit attribution and generalization. The
reference-implementation timings also leave open what optimized learned
heuristics and stronger large-scale classical controls would achieve.

\section{Conclusion}

The evaluated evolutionary policies do not establish a held-out advantage
over independent generation or a clear incremental benefit from the
prompt-feedback policy. Scoring diversity often produces trajectory
diversity without changing the final solution or cost. The form of
rarity weighting helps explain differences in frozen-program transfer,
while CP-SAT remains the stronger control on the original small graphs.
Broader representations and varied training corpora are needed to test
whether these search policies yield more useful deployed heuristics.

\section{Reproducibility}

Source code, configurations, selected heuristics and numerical results
are available in the accompanying repository. The README gives commands
for reproducing tables and analyses from saved observations without LLM
calls; \nolinkurl{scripts/analyze_selected.py} derives the winner-level
mechanism summaries. The original environment is Python 3.12.7 on
Windows 11, NetworkX 3.7, OR-Tools 9.15.6755, NumPy 2.5.4 and Matplotlib
3.11.2.

\subsection*{Generative AI disclosure}
Generative AI tools were used to support literature exploration, implementation, debugging, experimental analysis, and manuscript preparation. They also served as discussion aids for methodological decisions. The LLM used for heuristic generation is described separately in the methodology. The author reviewed the AI-assisted work and takes responsibility for the final results and manuscript.
